\begin{center}
  \Large
  \textbf{KATA PENGANTAR}
\end{center}

\addcontentsline{toc}{chapter}{KATA PENGANTAR}

\vspace{2ex}

Puji dan syukur penulis panjatkan ke hadirat Allah SWT, karena atas rahmat, karunia, dan ridha-Nya penulis dapat menyelesaikan penyusunan tugas akhir yang berjudul "\tatitle{}" dengan baik.

Dalam penyusunan tugas akhir ini, penulis menyadari bahwa banyak pihak yang telah memberikan bantuan, baik secara moril maupun materil, sehingga penulis dapat menyelesaikan tugas akhir ini. Oleh karena itu, pada kesempatan ini penulis ingin menyampaikan terima kasih yang sebesar-besarnya kepada:

\begin{enumerate}[nolistsep]
  \item Allah SWT yang senantiasa memberikan kesehatan, kemudahan, kekuatan, serta kelancaran kepada penulis dalam menyelesaikan tugas akhir ini.
  
  \item Kedua orang tua penulis yang selalu memberikan kasih sayang, dukungan, doa, serta motivasi sehingga penulis dapat menyelesaikan tugas akhir ini.
  
  \item Bapak Dr. Ir. Djoko Purwanto, M.Eng. dan Bapak Muhammad Qomaruz Zaman, S.T., M.T., Ph.D. selaku dosen pembimbing yang telah memberikan bimbingan, arahan, masukan, serta motivasi kepada penulis selama proses penyusunan tugas akhir ini.
  
  \item Seluruh dosen dan tenaga kependidikan di Departemen Teknik Elektro ITS yang telah memberikan ilmu pengetahuan, pengalaman, dan pelayanan yang baik selama penulis menempuh pendidikan.
  
  \item Raisyah Anna Sari Siregar yang selalu menjadi sumber semangat dan dukungan di setiap langkah penulisan ini, baik saat kita berjauhan maupun saat bersama, terima kasih untuk selalu percaya dan mengingatkanku untuk terus maju.
  
  \item Seluruh teman seperjuangan serta asisten Laboratorium MEST B202 yang selalu memberikan semangat, dukungan, bantuan, dan berbagai pengalaman berharga baik secara langsung maupun tidak langsung selama penyusunan tugas akhir ini.
  
  \item Seluruh pihak yang tidak dapat penulis sebutkan satu per satu yang telah membantu dan mendukung penulis dalam menyelesaikan tugas akhir ini.
\end{enumerate}

Penulis menyadari bahwa dalam penyusunan tugas akhir ini masih terdapat banyak kekurangan dan jauh dari kata sempurna. Oleh karena itu, penulis mengharapkan kritik dan saran yang membangun demi penyempurnaan penelitian ini di masa mendatang. Penulis berharap tugas akhir ini dapat memberikan manfaat serta menjadi referensi bagi pengembangan ilmu pengetahuan, khususnya pada bidang robotika, sistem navigasi otonom, dan \textit{computer vision}.

\begin{flushright}
  \begin{tabular}[b]{c}
    Surabaya, \MONTH{} \the\year{} \\
    \\
    \\
    \\
    \\
    Priagung Ramadhan Alfalakhi
  \end{tabular}
\end{flushright}